% Continuous editable transcription/translation of Volume I physical native leaves 23–26.
% Physical23 begins Article21 after completed Article20; no receiving sentence fragment.
% Physical26 ends a paragraph in Article29, which continues on physical27.
\SourceLeaf{23}

\Article{21}

\emph{Quando $a, b, c$ etc. sunt inter se primi, et productum $abc$ etc. potestas aliqua,
puta $=k^n$: singuli numeri $a, b, c$ etc. similes potestates erunt.}

Sit $a = l^\lambda m^\mu p^\pi$ etc., designantibus $l, m, p$ etc. numeros primos diversos,
quorum nullus per \emph{hyp.} est factor numerorum $b, c$ etc. Quare productum $abc$ etc.
factorem $l$ implicabit $\lambda$ vicibus, factorem $m$ vero $\mu$ vicibus etc.: hinc
(\emph{art. praec.}) $\lambda, \mu, \pi$ etc. per $n$ divisibiles adeoque
\[
{}^{n}\!\surd a = l^{\frac{\lambda}{n}} \, m^{\frac{\mu}{n}} \, p^{\frac{\pi}{n}} \text{ etc.}
\]
integer. Similiter de reliquis $b, c$ etc.

Haec de numeris primis praemittenda erant; iam ad ea quae finem nobis propositum propius
attinent convertimur.

\Article{22}

\emph{Si numeri $a, b$ per alium $k$ divisibiles secundum modulum $m$ ad $k$ primum sunt
congrui: $\frac{a}{k}$ et $\frac{b}{k}$ secundum eundem modulum congrui erunt.}

Patet enim $a-b$ per $k$ divisibilem fore, nec minus per $m$ (\emph{hyp.}); quare
(\emph{art. 19}) $\frac{a-b}{k}$ per $m$ divisibilis erit, i. e. erit
$\frac{a}{k} \equiv \frac{b}{k}$ (mod. $m$).

Si autem reliquis manentibus $m$ et $k$ habent divisorem communem maximum $e$, erit
$\frac{a}{k} \equiv \frac{b}{k}$ (mod. $\frac{m}{e}$). Namque $\frac{k}{e}$ et
$\frac{m}{e}$ inter se primi. At $a-b$ tam per $k$ quam per $m$ divisibilis adeoque etiam
$\frac{a-b}{e}$ tam per $\frac{k}{e}$ quam per $\frac{m}{e}$, hincque per $\frac{km}{ee}$
i. e. $\frac{a-b}{k}$ per $\frac{m}{e}$, sive $\frac{a}{k} \equiv \frac{b}{k}$ (mod.
$\frac{m}{e}$).

\Article{23}

\emph{Si $a$ ad $m$ primus, et $e, f$ numeri secundum modulum $m$ incongrui: erunt etiam
$ae, af$ incongrui secundum $m$.}

Hoc est tantum conversio theor. art. praec.

Hinc vero manifestum est, si $a$ per omnes numeros integros a $0$ usque ad $m-1$
multiplicetur productaque secundum modulum $m$ ad residua sua minima reducantur, haec omnia
fore inaequalia. Et quum horum residuorum, quorum nullum $> m$, numerus sit $m$, totidemque
dentur numeri a $0$ usque ad $m-1$, patet, nullum horum numerorum inter illa residua deesse
posse

\SourceLeaf{24}

\Article{24}

\emph{Expressio $ax+b$}, denotantibus $a, b$ numeros datos, $x$ numerum indeterminatum seu
variabilem, \emph{secundum modulum $m$, ad $a$ primum, cuivis numero dato congrua fieri potest.}

Sit numerus, cui congrua fieri debet, $c$, et residuum minimum positivum ipsius $c-b$
secundum modulum $m$, $e$. Ex art. praec. necessario datur valor ipsius $x < m$, talis, ut
producti $ax$ secundum modulum $m$ residuum minimum fiat $e$; esto hic valor $v$, eritque
$av \equiv e \equiv c-b$; unde $av+b \equiv c$ (mod. $m$) \emph{Q. E. F.}

\Article{25}

Expressionem duas quantitates congruas exhibentem ad instar aequationum,
\emph{congruentiam} vocamus; quae si incognitam implicat, \emph{resolvi} dicitur, quando
pro hac valor invenitur congruentiae satisfaciens (\emph{radix}). Hinc porro intelligitur,
quid sit \emph{congruentia resolubilis et congruentia irresolubilis}. Tandem facile
perspicitur similes distinctiones locum hic habere posse uti in aequationibus.
Congruentiarum \emph{transscendentium} 
infra exempla occurrent: \emph{algebraicae} vero secundum dimensionem maximam
incognitaè 
in congruentias primi, secundi altiorumque \emph{graduum} distribuuntur. Nec minus
congruentiae plures proponi possunt plures incognitas involventes, de quarum
\emph{eliminatione} disquirendum.

\begin{center}
\emph{Solutio congruentiarum primi gradus.}
\end{center}

\Article{26}

Congruentia itaque primi gradus $ax+b \equiv c$ ex art. 24 semper resolubilis, quando
modulus ad $a$ est primus. Quodsi vero $v$ fuerit valor idoneus ipsius $x$, sive radix
congruentiae, palam est, omnes numeros, ipsi $v$ secundum congruentiae propositae modulum
congruos, etiam radices fore (art. 9). Neque minus facile perspicitur, omnes radices ipsi
$v$ congruos esse debere: si enim alia radix fuerit $t$, erit $av+b \equiv at+b$, unde
$av \equiv at$, et hinc $v \equiv t$ (art. 22). Hinc colligitur congruentiam
$x \equiv v$ (mod. $m$) exhibere resolutionem completam congruentiae $ax+b \equiv c$.

Quia resolutiones congruentiae per valores ipsius $x$ congruos per se sunt obviae, atque,
hoc respectu, numeri congrui tamquam aequivalentes considerandi, tales congruentiae
resolutiones pro una eademque habebimus. Quamobrem quum
\SourceLeaf{25}
nostra congruentia $ax+b \equiv c$ alias resolutiones non admittat, pronunciabimus, unico
tantum modo eam esse resolubilem seu unam tantum radicem habere. Ita \emph{e. g.}
congruentia $6x+5 \equiv 13$ (mod. $11$) alias radices non admittit, quam quae sunt
$\equiv 5$ (mod. $11$). Haud perinde res se habet in congruentiis altiorum graduum, sive
etiam in congruentiis primi gradus, ubi incognita per numerum est multiplicata, ad quem
modulus non est primus.

\Article{27}

Superest, ut de invenienda resolutione ipsa congruentiae huiusmodi, quaedam adiiciamus.
Primo observamus, congruentiam formae $ax+t \equiv u$, cuius modulum ad $a$ primum
supponimus, ab hac $ax \equiv \pm 1$ pendere: si enim huic satisfacit $x \equiv r$, illi
satisfaciet $x \equiv \pm(u-t)r$. At congruentiae $ax \equiv \pm 1$, modulo per $b$
designato, aequivalet aequatio indeterminata $ax = by \pm 1$, quae quomodo sit solvenda
hoc quidem tempore abunde est notum; quare nobis sufficiet, calculi algorithmum huc
transscripsisse.

Si quantitates $A, B, C, D, E$ etc. ita ab his $\alpha, \beta, \gamma, \delta$ etc.
pendent, ut habeatur
\[
A = \alpha, \quad B = \beta A + 1, \quad C = \gamma B + A, \quad D = \delta C + B, \quad
E = \varepsilon D + C \text{ etc.}
\]
brevitatis gratia ita eas designamus,
\begingroup
\renewcommand{\thefootnote}{*)}
\[
A = [\alpha], \quad B = [\alpha, \beta], \quad C = [\alpha, \beta, \gamma], \quad
D = [\alpha, \beta, \gamma, \delta] \text{ etc.}\footnotemark.
\]
\footnotetext{
Multo generalius haecce relatio considerari potest, quod negotium alia forsan occasione
suscipiemus. Hic duas tantum propositiones adiicimus, quae usum suum in praesenti
investigatione habent; scilicet,

\[1^{\mathrm{o}}.\quad [\alpha,\beta,\gamma\mathinner{\ldotp\ldotp\ldotp\ldotp} \lambda,\mu].[\beta,\gamma\mathinner{\ldotp\ldotp\ldotp\ldotp} \lambda] -
[\alpha,\beta,\gamma\ldots \lambda][\beta,\gamma\ldots \lambda,\mu] = \pm 1,\]
ubi signum superius accipiendum quando numerorum $\alpha,\beta,\gamma\ldots \lambda,\mu$
multitudo par, inferius quando impar.

$2^{\mathrm{o}}$. Numerorum $\alpha,\beta,\gamma$ etc. ordo inverti potest,
$[\alpha,\beta,\gamma\ldots \lambda,\mu]=[\mu,\lambda\mathinner{\ldotp\ldotp\ldotp\ldotp}\gamma,\beta,\alpha].$

Demonstrationes quae non sunt difficiles hic supprimimus.
}
\endgroup
Iam proposita sit aequatio indeterminata $ax = by \pm 1$, ubi $a, b$ positivi. Supponamus,
id quod licet, $a$ esse non $< b$. Tum ad instar algorithmi noti, secundum quem duorum
numerorum divisor communis maximus investigatur, formentur per divisionem vulgarem
aequationes,
\[
a = \alpha b + c, \quad b = \beta c + d, \quad c = \gamma d + e \text{ etc.}
\]
ita ut $\alpha, \beta, \gamma$ etc. $c, d, e$ etc. sint integri positivi, et $b, c, d, e$
continuo decrescentes, donec perveniatur ad $m = \mu n + 1$
\SourceLeaf{26}
quod tandem evenire debere constat. Erit itaque
\[
a=[n,\mu,\mathinner{\ldotp\ldotp\ldotp\ldotp} \gamma,\beta,\alpha], \quad b=[n,\mu\mathinner{\ldotp\ldotp\ldotp\ldotp} \gamma,\beta]
\]
Tum fiat
\[
x=[\mu,\mathinner{\ldotp\ldotp\ldotp\ldotp} \gamma,\beta], \quad y=[\mu,\mathinner{\ldotp\ldotp\ldotp\ldotp} \gamma,\beta,\alpha]
\]
eritque $ax=by+1$, quando numerorum $\alpha,\beta,\gamma\mathinner{\ldotp\ldotp\ldotp\ldotp} \mu,n$ multitudo est par, aut
$ax=by-1$, quando est impar. \emph{Q. E. F.}

\Article{28}

Resolutionem generalem huiusmodi aequationum indeterminatarum ill. Euler primus docuit,
\emph{Comment. Petrop. T.} VII. \emph{p.} 46. Methodus qua usus est consistit in substitutione
aliarum incognitarum loco ipsarum $x, y$, atque hoc quidem tempore satis est nota. Ill. La
Grange paullo aliter rem aggressus est: scilicet ex theoria fractionum continuarum constat,
si fractio $\frac{b}{a}$ in fractionem continuam
\[
\cfrac{1}{\alpha + \cfrac{1}{\beta + \cfrac{1}{\gamma + \text{etc.} + \cfrac{1}{\mu + \frac{x}{n}}}}}
\]
convertatur, haecque deleta ultima sui parte $\frac{x}{n}$ in fractionem communem
$\frac{x}{y}$ restituatur, fore $ax=by\pm1$, siquidem fuerit $a$ ad $b$ primus. Ceterum ex
utraque methodo idem algorithmus derivatur. Investigationes ill. La Grange exstant
\emph{Hist. de l'Ac. de Berlin Année} 1767 \emph{p.} 175, et cum aliis in \emph{Supplementis
versioni gallicae Algebrae Eulerianae adiectis}.

\Article{29}

Congruentiae $ax+t \equiv u$ cuius modulus ad $a$ non primus, facile ad casum praecedentem
reducitur. Sit modulus $m$, maximusque numerorum $a, m$ divisor communis $\delta$. Primo
patet quemvis valorem ipsius $x$ congruentiae secundum modulum $m$ satisfacientem eidem
etiam secundum modulum $\delta$ satisfacere (art. 5). At semper $ax \equiv 0$ (mod.
$\delta$), quoniam $\delta$ ipsum $a$ metitur. Quare, nisi $t \equiv u$ (mod. $\delta$) i.
e. $t-u$ per $\delta$ divisibilis, congruentia proposita non est resolubilis.
% End of validated physical26 scope; Article29 and Section II remain open.
